\section{Resultados}

\subsection{Efeito da taxonomia}

A Tabela~\ref{tab:main} resume os modelos treinados. Material elevou o AP$_{50}$ médio de 0,146 para 0,183 e o $F_1$ de 0,232 para 0,296 em relação a Fine. Binary alcançou 0,567 de AP$_{50}$ e apresentou a menor variação entre sementes ($\sigma=0{,}002$). Esse ganho não constitui uma comparação de capacidade semântica equivalente: Binary só precisa separar resíduo de fundo.

\begin{table}[H]
  \centering
  \caption{Teste por taxonomia: média $\pm$ desvio-padrão de três sementes.}
  \label{tab:main}
  \resizebox{\linewidth}{!}{\begin{tabular}{lrrrrr}
\toprule
Taxonomia & $P$ & $R$ & $F_1$ & AP$_{50}$ & AP$_{50:95}$ \\
\midrule
Fine & 0.360$\pm$0.063 & 0.173$\pm$0.010 & 0.232$\pm$0.005 & 0.146$\pm$0.022 & 0.118$\pm$0.022 \\
Material & 0.387$\pm$0.050 & 0.241$\pm$0.012 & 0.296$\pm$0.006 & 0.183$\pm$0.023 & 0.131$\pm$0.023 \\
Binary & 0.750$\pm$0.050 & 0.491$\pm$0.018 & 0.593$\pm$0.004 & 0.567$\pm$0.002 & 0.405$\pm$0.005 \\
\bottomrule
\end{tabular}
}
\end{table}

O bootstrap separa média superior de evidência repetível. Embora Material tenha superado Fine em média nas três métricas, seus ganhos de AP tiveram IC95\% acima de zero em apenas uma das três sementes. O ganho de $F_1$, por outro lado, foi sustentado nas três. Portanto, os modelos treinados separadamente permitem afirmar uma melhora consistente em $F_1$, mas não uma superioridade inequívoca de AP.

\subsection{Localização versus semântica}

A Tabela~\ref{tab:hierarchy} quantifica a variação observada quando as \emph{mesmas} caixas e confianças produzidas pelo modelo Fine são avaliadas com menor granularidade semântica. Ao aceitar a macroclasse de material, AP$_{50}$ passou de $0{,}103$ para $0{,}171$; sem distinção entre classes, alcançou $0{,}466$. Como esta análise não envolve novo treinamento, o ganho estima a parcela da penalização associada aos rótulos, embora a alteração da macro-média impeça sua interpretação como contagem pura de erros semânticos. A Tabela~\ref{tab:bootstrap} apresenta a repetibilidade desse efeito entre sementes.

\begin{table}[H]
  \centering
  \caption{Mesmas predições Fine avaliadas em três níveis pelo avaliador controlado (média $\pm$ desvio-padrão).}
  \label{tab:hierarchy}
  \begin{tabular}{lrrr}
\toprule
Nível de avaliação & $F_1$ & AP$_{50}$ & AP$_{50:95}$ \\
\midrule
Fine & 0.168$\pm$0.016 & 0.103$\pm$0.007 & 0.082$\pm$0.005 \\
Material & 0.264$\pm$0.021 & 0.171$\pm$0.015 & 0.122$\pm$0.010 \\
Binary & 0.506$\pm$0.011 & 0.466$\pm$0.014 & 0.322$\pm$0.011 \\
\bottomrule
\end{tabular}

\end{table}

\begin{table}[H]
  \centering
  \footnotesize
  \caption{Síntese do bootstrap pareado nas três sementes.}
  \label{tab:bootstrap}
  \begin{tabular}{llrr}
\toprule
Comparação ($A-B$) & Métrica & $\overline{\Delta}\pm\sigma$ & IC exclui 0 \\
\midrule
Fine$-$Material & AP$_{50}$ & $-0{,}051\pm0{,}022$ & 1/3 \\
Fine$-$Material & AP$_{50:95}$ & $-0{,}029\pm0{,}023$ & 1/3 \\
Fine$-$Material & $F_1$ & $-0{,}074\pm0{,}027$ & 3/3 \\
Fine$\rightarrow$Material$-$Fine & AP$_{50}$ & $+0{,}068\pm0{,}012$ & 3/3 \\
Fine$\rightarrow$Material$-$Fine & AP$_{50:95}$ & $+0{,}039\pm0{,}010$ & 2/3 \\
Fine$\rightarrow$Material$-$Fine & $F_1$ & $+0{,}101\pm0{,}008$ & 3/3 \\
Fine$\rightarrow$Binary$-$Fine & AP$_{50}$ & $+0{,}345\pm0{,}016$ & 3/3 \\
Fine$\rightarrow$Binary$-$Fine & AP$_{50:95}$ & $+0{,}226\pm0{,}012$ & 3/3 \\
Fine$\rightarrow$Binary$-$Fine & $F_1$ & $+0{,}318\pm0{,}011$ & 3/3 \\
\bottomrule
\end{tabular}

\end{table}

Cada célula da Tabela~\ref{tab:bootstrap} apresenta ganho médio $\pm$ desvio-padrão e, entre parênteses, quantos ICs de 95\% excluem zero. Valores positivos favorecem a alternativa mais simples; 3/3 indica evidência repetida nas três sementes, enquanto 1/3 recomenda cautela. ``Fine como'' denota remapeamento das mesmas predições, sem novo treinamento.

\begin{figure}[H]
  \centering
  \includegraphics[width=\linewidth]{taxonomy_and_hierarchy.pdf}
  \caption{À esquerda, modelos treinados separadamente e avaliados pelo Ultralytics; à direita, mesmas predições Fine remapeadas e avaliadas pelo protocolo controlado. Barras de erro: um desvio-padrão entre sementes.}
  \label{fig:hierarchy}
\end{figure}

\subsection{Exemplos nos níveis Material e Binary}

A Figura~\ref{fig:model-level} amplia a análise visual para os demais modelos. Em (a), o modelo Material localiza e classifica corretamente um filme plástico. Os painéis (b) e (c) utilizam a mesma garrafa: o modelo Material preserva a localização, mas prediz Paper/Cardboard em vez de Plastic, enquanto o modelo Binary a identifica como \textit{Litter}. A comparação evidencia que a remoção da exigência de material pode converter uma confusão semântica em detecção correta. O painel (d) delimita esse efeito: apesar da ausência de distinção entre materiais, o modelo Binary omite um pneu de grandes dimensões no limiar de confiança 0,25. A redução da granularidade elimina determinados erros semânticos, mas não assegura cobertura espacial.

\begin{figure}[H]
  \centering
  \includegraphics[width=.94\linewidth]{model_level_examples.pdf}
  \caption{Predições reais dos modelos Material e Binary, semente 42. Os painéis (b) e (c) mostram a mesma imagem sob duas taxonomias; (d) mostra uma omissão binária.}
  \label{fig:model-level}
\end{figure}

\subsection{Cauda longa e tipos de erro}

A Tabela~\ref{tab:tail} apresenta crescimento monotônico do desempenho com a frequência. Entre as classes presentes no teste, as raras obtiveram AP$_{50}=0{,}046$ e revocação $0{,}053$, enquanto as frequentes atingiram 0,207 e 0,282. Esse resultado indica que a dificuldade da taxonomia Fine não se distribui uniformemente entre os rótulos.

\begin{table}[H]
  \centering
  \caption{Desempenho Fine por frequência no treino, agregado em três sementes; somente classes com suporte no teste.}
  \label{tab:tail}
  \begin{tabular}{lrrrr}
\toprule
Grupo & \#classes & $n_{treino}$ & AP$_{50}$ & Recall \\
\midrule
Raras & 28 & 5.1 & 0.046$\pm$0.009 & 0.030$\pm$0.007 \\
Médias & 22 & 43.9 & 0.098$\pm$0.022 & 0.198$\pm$0.049 \\
Frequentes & 10 & 225.1 & 0.207$\pm$0.017 & 0.282$\pm$0.019 \\
\bottomrule
\end{tabular}

\end{table}

Como diagnóstico adicional, no limiar de confiança 0,25, as três sementes produziram $100{,}0\pm5{,}2$ acertos Fine e $132{,}3\pm7{,}2$ confusões semânticas entre $232{,}3\pm10{,}6$ casamentos espaciais. Portanto, $56{,}95\pm1{,}38\%$ dos casamentos em IoU$\geq0{,}5$ tinham classe Fine incorreta. Em média, houve ainda $16{,}3\pm5{,}5$ falhas de localização ($0{,}1\leq$IoU$<0{,}5$), $40{,}7\pm8{,}1$ duplicatas, $47{,}3\pm5{,}0$ falsos positivos de fundo (IoU$<0{,}1$) e $404{,}7\pm10{,}6$ objetos omitidos. Estes correspondem a $63{,}5\pm1{,}7\%$ dos 637 objetos de teste. Os totais dependem do limiar e servem para caracterização, não como substitutos das curvas de AP \cite{bolya2020tide}.

\begin{table}[H]
  \centering
  \small
  \caption{Decomposição dos erros Fine por semente no limiar de confiança 0,25. Categorias de predição e omissões são mutuamente exclusivas dentro de seus respectivos inventários.}
  \label{tab:errors}
  \begin{tabular}{lrrrr}
\toprule
Categoria & 42 & 123 & 2026 & Média $\pm$ DP \\
\midrule
Acerto Fine & 97 & 106 & 97 & $100{,}0\pm5{,}2$ \\
Confusão no material & 51 & 57 & 52 & $53{,}3\pm3{,}2$ \\
Confusão entre materiais & 73 & 79 & 85 & $79{,}0\pm6{,}0$ \\
Falha de localização & 11 & 16 & 22 & $16{,}3\pm5{,}5$ \\
Duplicata & 32 & 42 & 48 & $40{,}7\pm8{,}1$ \\
Falso positivo de fundo & 42 & 48 & 52 & $47{,}3\pm5{,}0$ \\
Objeto omitido & 416 & 395 & 403 & $404{,}7\pm10{,}6$ \\
\bottomrule
\end{tabular}

\end{table}

A Figura~\ref{fig:qualitative} liga essa decomposição às imagens de teste. A primeira linha apresenta o mesmo recorte sem anotações, com a verdade-terreno e com uma saída correta. Na segunda, uma lata é localizada corretamente, mas confundida com aerossol, preservando o material; uma embalagem metálica é confundida com papel/cartão; e uma tampa plástica anotada não recebe predição. Os exemplos tornam visíveis erros semânticos e espaciais de naturezas distintas.

\begin{figure}[!t]
  \centering
  \includegraphics[width=.94\linewidth]{qualitative_analysis.pdf}
  \caption{Casos reais do teste Fine, semente 42. A primeira linha separa entrada, anotação e saída correta; a segunda exemplifica confusão intra-material, confusão inter-material e omissão. As imagens foram recortadas apenas para legibilidade; a prevalência de cada erro é quantificada nas três sementes na Tabela~\ref{tab:errors}.}
  \label{fig:qualitative}
\end{figure}
